%=============================================================================
% Chapter 2: Advanced Function Concepts
% Structured Programming with C++ -- Lecture 12
%=============================================================================

\chapter{Advanced Function Concepts}
\label{ch:advanced-functions}

%-----------------------------------------------------------------------------
\section{Types of User-Defined Functions}
\label{sec:function-types}
%-----------------------------------------------------------------------------

User-defined functions in C++ can be classified into three categories based on
whether they accept arguments and whether they return a value.

\begin{infobox}{Classification of Functions}
\begin{enumerate}
    \item \textbf{No arguments, no return value} --- The function neither
          receives data from the caller nor sends data back. It is declared
          as \texttt{void functionName(void)} or
          \texttt{void functionName()}.
    \item \textbf{With arguments, no return value} --- The function receives
          data through its parameters but does not return a result. It is
          declared as \texttt{void functionName(parameters)}.
    \item \textbf{With arguments, with return value} --- The function
          receives data and returns a computed result. It is declared as
          \texttt{type functionName(parameters)}.
\end{enumerate}
\end{infobox}

%-----------------------------------------------------------------------------
\subsection{Function Types Summary Table}
%-----------------------------------------------------------------------------

\begin{center}
\renewcommand{\arraystretch}{1.4}
\footnotesize
\begin{tabular}{@{}cllp{3cm}p{3.2cm}@{}}
\toprule
\textbf{Type} & \textbf{Arguments} & \textbf{Return}
    & \textbf{Example Signature}
    & \textbf{Use Case} \\
\midrule
1 & None & None
    & \texttt{void greet()}
    & Print a menu / banner \\
2 & Yes  & None
    & \texttt{void printSquare(int n)}
    & Display formatted output \\
3 & Yes  & Yes
    & \texttt{float calcSquare(float n)}
    & Compute \& return a result \\
\bottomrule
\end{tabular}
\end{center}

%-----------------------------------------------------------------------------
\subsection{Type 1: No Arguments, No Return Value}
%-----------------------------------------------------------------------------

\begin{cppcode}
#include <iostream>
using namespace std;

void greet()
{
    cout << "Hello! Welcome to C++ programming."
         << endl;
    cout << "Let's learn about functions!" << endl;
}

int main()
{
    greet();   // No arguments, no return value
    return 0;
}
\end{cppcode}

\noindent\textbf{Output:}
\begin{verbatim}
Hello! Welcome to C++ programming.
Let's learn about functions!
\end{verbatim}

%-----------------------------------------------------------------------------
\subsection{Type 2: With Arguments, No Return Value}
%-----------------------------------------------------------------------------

\begin{cppcode}
#include <iostream>
using namespace std;

void printSquare(int n)
{
    cout << "The square of " << n
         << " is: " << n * n << endl;
}

int main()
{
    int x = 7;
    printSquare(x);   // Argument passed, no return
    printSquare(12);  // Can also pass a literal
    return 0;
}
\end{cppcode}

\noindent\textbf{Output:}
\begin{verbatim}
The square of 7 is: 49
The square of 12 is: 144
\end{verbatim}

%-----------------------------------------------------------------------------
\subsection{Type 3: With Arguments, With Return Value}
%-----------------------------------------------------------------------------

\begin{cppcode}
#include <iostream>
using namespace std;

float calcSquare(float n)
{
    return n * n;
}

int main()
{
    float x = 3.5;
    float result = calcSquare(x);
    cout << "The square of " << x
         << " is: " << result << endl;

    // Returned value used directly in an expression
    cout << "Twice the square: "
         << 2 * calcSquare(x) << endl;
    return 0;
}
\end{cppcode}

\noindent\textbf{Output:}
\begin{verbatim}
The square of 3.5 is: 12.25
Twice the square: 24.5
\end{verbatim}

\begin{conceptbox}{Choosing the Right Type}
\begin{itemize}
    \item Use \textbf{Type 1} for standalone actions such as printing a menu
          or a welcome message.
    \item Use \textbf{Type 2} when the function needs input but only performs
          an action (e.g., displaying output) without producing a result.
    \item Use \textbf{Type 3} when the function computes and returns a value
          that the caller needs for further processing.
\end{itemize}
\end{conceptbox}

%-----------------------------------------------------------------------------
\subsection{Side-by-Side Comparison: \texttt{void} vs.\ Return Value}
%-----------------------------------------------------------------------------

Consider two approaches to computing the square of a number:

\begin{cppcode}
#include <iostream>
using namespace std;

// Approach A: void -- prints directly
void squareVoid(float n)
{
    cout << "Square (void): " << n * n << endl;
}

// Approach B: returns the result
float squareReturn(float n)
{
    return n * n;
}

int main()
{
    float num = 4.0;

    // Approach A: printed inside the function
    squareVoid(num);

    // Approach B: returned and used further
    float res = squareReturn(num);
    cout << "Square (return): " << res << endl;

    // Returned value in an expression
    cout << "Twice the square: "
         << 2 * squareReturn(num) << endl;

    return 0;
}
\end{cppcode}

\noindent\textbf{Output:}
\begin{verbatim}
Square (void): 16
Square (return): 16
Twice the square: 32
\end{verbatim}

\begin{importantbox}{Return Values Enable Composition}
Functions that \emph{return} values are more flexible because their results
can be stored in variables, passed to other functions, or used in expressions.
The \texttt{void} approach is suitable only when the sole purpose of the
function is to produce a side effect (such as printing).
\end{importantbox}

%-----------------------------------------------------------------------------
\section{Actual and Formal Arguments}
\label{sec:actual-formal-args}
%-----------------------------------------------------------------------------

When working with functions, it is essential to distinguish between
\textbf{actual arguments} and \textbf{formal arguments}.

\begin{definitionbox}{Actual Arguments (Actual Parameters)}
\textbf{Actual arguments} are the variables, constants, or expressions
specified in the \emph{function call}. They represent the real data that is
passed to the function when it is invoked.
\end{definitionbox}

\begin{definitionbox}{Formal Arguments (Formal Parameters)}
\textbf{Formal arguments} (also called \emph{dummy variables} or
\emph{parametric variables}) are the variables declared in the
\emph{function definition} header. They act as placeholders that receive the
values of the actual arguments when the function is called.
\end{definitionbox}

The following example illustrates the distinction clearly:

\begin{cppcode}
#include <iostream>
using namespace std;

// Function prototype
int add(int a, int b);  // a, b = formal arguments

int main()
{
    int x = 5, y = 3;

    // x and y are ACTUAL arguments
    int sum = add(x, y);
    cout << "Sum = " << sum << endl;

    // Constants can also be actual arguments
    cout << "Sum = " << add(10, 20) << endl;

    return 0;
}

// Function definition
int add(int a, int b)   // a, b = FORMAL arguments
{
    // a receives value of x, b receives value of y
    return a + b;
}
\end{cppcode}

\noindent\textbf{Output:}
\begin{verbatim}
Sum = 8
Sum = 30
\end{verbatim}

\begin{conceptbox}{Mapping Between Actual and Formal Arguments}
When \texttt{add(x, y)} is called:
\begin{itemize}
    \item The actual argument \texttt{x} (value 5) is copied into the formal
          argument \texttt{a}.
    \item The actual argument \texttt{y} (value 3) is copied into the formal
          argument \texttt{b}.
\end{itemize}
The names of actual and formal arguments need not match. What matters is the
\textbf{order}, \textbf{number}, and \textbf{type} of the arguments ---
these must be consistent between the call and the definition.
\end{conceptbox}

\begin{center}
\renewcommand{\arraystretch}{1.4}
\begin{tabular}{@{}lll@{}}
\toprule
\textbf{Property} & \textbf{Actual Arguments}
    & \textbf{Formal Arguments} \\
\midrule
Appear in          & Function \emph{call}
    & Function \emph{definition} \\
Also called        & Real parameters
    & Dummy / parametric variables \\
Examples           & \texttt{x}, \texttt{y}, \texttt{10}
    & \texttt{a}, \texttt{b} \\
Matching rule      & \multicolumn{2}{l}{%
    Must agree in order, number, and type} \\
\bottomrule
\end{tabular}
\end{center}

\begin{warningbox}{Common Mistakes with Arguments}
\begin{itemize}
    \item \textbf{Mismatched types:} Passing a \texttt{float} where an
          \texttt{int} is expected may cause implicit truncation.
    \item \textbf{Wrong number of arguments:} The number of actual arguments
          must match the number of formal parameters (unless default values
          are provided).
    \item \textbf{Wrong order:} Arguments are matched by \emph{position}, not
          by name. Swapping the order can produce incorrect results.
\end{itemize}
\end{warningbox}

%-----------------------------------------------------------------------------
\section{Local and Global Variables}
\label{sec:local-global-variables}
%-----------------------------------------------------------------------------

The \textbf{scope} of a variable determines where in a program it can be
accessed. C++ distinguishes between local and global scope.

%- - - - - - - - - - - - - - - - - - - - - - - - - - - - - - - - - - - - - -
\subsection{Local Variables}
\label{subsec:local-variables}
%- - - - - - - - - - - - - - - - - - - - - - - - - - - - - - - - - - - - - -

\begin{definitionbox}{Local Variable}
A \textbf{local variable} is a variable defined \emph{inside} a function or a
block (enclosed in \texttt{\{\}}). It exists only during the execution of that
block and is destroyed when the block ends. Local variables cannot be accessed
from outside their enclosing block.
\end{definitionbox}

\begin{cppcode}
#include <iostream>
using namespace std;

void func(int i, int j)
{
    int k, m;         // k and m are LOCAL to func()
    k = i + j;
    m = i * j;
    cout << "k = " << k << ", m = " << m << endl;
}

int main()
{
    int a = 4, b = 5; // a and b are LOCAL to main()
    func(a, b);

    // cout << k;  // ERROR: k is not accessible here
    return 0;
}
\end{cppcode}

\noindent\textbf{Output:}
\begin{verbatim}
k = 9, m = 20
\end{verbatim}

\begin{infobox}{Key Properties of Local Variables}
\begin{itemize}
    \item Local variables \textbf{belong to the block} in which they are
          declared.
    \item Variables with the \textbf{same name} declared in different
          functions (or blocks) are \textbf{completely independent entities},
          each with its own memory location.
    \item Local variables are allocated on the \textbf{stack} and are
          automatically destroyed when the function returns.
    \item Function parameters (\texttt{i} and \texttt{j} above) are also
          treated as local variables within the function.
\end{itemize}
\end{infobox}

%- - - - - - - - - - - - - - - - - - - - - - - - - - - - - - - - - - - - - -
\subsection{Global Variables}
\label{subsec:global-variables}
%- - - - - - - - - - - - - - - - - - - - - - - - - - - - - - - - - - - - - -

\begin{definitionbox}{Global Variable}
A \textbf{global variable} is a variable defined \emph{outside} all functions,
typically at the top of the source file. It retains the same data type and
name throughout the entire program and can be accessed from any function.
\end{definitionbox}

\begin{cppcode}
#include <iostream>
using namespace std;

// Global variables -- accessible from any function
int x;
int y = 5;

void sum(void)
{
    // Reads and modifies global variables
    int s = x + y;
    cout << "x = " << x << ", y = " << y << endl;
    cout << "Sum = " << s << endl;
}

int main()
{
    x = 10;      // Assign to global x
    sum();

    x = 100;
    y = 200;
    sum();        // Globals have new values now

    return 0;
}
\end{cppcode}

\noindent\textbf{Output:}
\begin{verbatim}
x = 10, y = 5
Sum = 15
x = 100, y = 200
Sum = 300
\end{verbatim}

\begin{importantbox}{Global Variable Initialization}
Global variables in C++ are automatically \textbf{zero-initialized} if no
explicit initial value is provided. In the example above, \texttt{int x;}
starts at~0, while \texttt{int y = 5;} starts at~5.
\end{importantbox}

Here is a second example showing how multiple functions share global data:

\begin{cppcode}
#include <iostream>
using namespace std;

// Global variables
int x, y;

void readValues()
{
    cout << "Enter two integers: ";
    cin >> x >> y;    // Modifying global variables
}

void printSum()
{
    int sum = x + y;  // Reading global variables
    cout << "The sum of " << x << " and " << y
         << " is: " << sum << endl;
}

int main()
{
    readValues();
    printSum();
    return 0;
}
\end{cppcode}

\noindent\textbf{Sample Output:}
\begin{verbatim}
Enter two integers: 12 8
The sum of 12 and 8 is: 20
\end{verbatim}

\begin{warningbox}{Use Global Variables Sparingly}
While global variables provide a convenient way to share data between
functions, they have significant drawbacks:
\begin{itemize}
    \item Any function can modify a global variable, making it difficult to
          track where changes occur.
    \item Global variables increase \textbf{coupling} between functions,
          reducing modularity.
    \item They can lead to subtle bugs when two functions unintentionally
          modify the same variable.
\end{itemize}
Prefer passing data through \textbf{parameters and return values} whenever
possible. Reserve global variables for data that genuinely needs to be
shared across many functions.
\end{warningbox}

\begin{conceptbox}{Scope Summary}
\begin{center}
\renewcommand{\arraystretch}{1.3}
\begin{tabular}{@{}llll@{}}
\toprule
\textbf{Property} & \textbf{Local Variable}
    & \textbf{Global Variable} \\
\midrule
Declared in      & Inside a function/block
    & Outside all functions \\
Accessible from  & Only its own block
    & Any function \\
Lifetime         & Block execution
    & Entire program \\
Default value    & Undefined (garbage)
    & Zero-initialized \\
Storage          & Stack
    & Data segment \\
\bottomrule
\end{tabular}
\end{center}
\end{conceptbox}

%-----------------------------------------------------------------------------
\section{Recursive Functions}
\label{sec:recursive-functions}
%-----------------------------------------------------------------------------

\begin{definitionbox}{Recursion}
\textbf{Recursion} is a programming technique in which a function calls
\emph{itself}, either directly or indirectly. A function that employs
recursion is called a \textbf{recursive function}.
\end{definitionbox}

Recursion is a powerful tool that is particularly useful for problems that can
be broken down into smaller instances of the same problem, such as traversing
\textbf{linked lists}, searching \textbf{trees}, and computing mathematical
sequences.

\begin{importantbox}{Two Essential Components of a Recursive Function}
Every recursive function must have:
\begin{enumerate}
    \item \textbf{Base case:} A condition under which the function does
          \emph{not} call itself, preventing infinite recursion.
    \item \textbf{Recursive case:} The function calls itself with a modified
          argument that moves toward the base case.
\end{enumerate}
Without a proper base case, the function will call itself indefinitely,
eventually causing a \textbf{stack overflow} error.
\end{importantbox}

\begin{infobox}{Recursion vs.\ Iteration}
\begin{itemize}
    \item A \textbf{normal (iterative) function} uses loops
          (\texttt{for}, \texttt{while}) to repeat operations.
    \item A \textbf{recursive function} achieves repetition by calling itself
          with progressively simpler inputs.
    \item Any problem solvable with recursion can also be solved iteratively,
          and vice versa. The choice depends on clarity and efficiency.
\end{itemize}
\end{infobox}

%-----------------------------------------------------------------------------
\subsection{Example: Recursive Sum
    \texorpdfstring{$1 + 2 + \cdots + n$}{1+2+...+n}}
\label{subsec:recursive-sum}
%-----------------------------------------------------------------------------

The sum of the first $n$ positive integers can be defined recursively as:
\[
    \text{sum}(n) =
    \begin{cases}
        0 & \text{if } n = 0 \quad \text{(base case)} \\
        n + \text{sum}(n - 1) & \text{if } n > 0
            \quad \text{(recursive case)}
    \end{cases}
\]

\begin{cppcode}
#include <iostream>
using namespace std;

int sum(int n)
{
    if (n == 0)
        return 0;           // Base case
    else
        return n + sum(n - 1);  // Recursive case
}

int main()
{
    int n;
    cout << "Enter a positive integer: ";
    cin >> n;

    cout << "Sum from 1 to " << n
         << " = " << sum(n) << endl;

    return 0;
}
\end{cppcode}

\noindent\textbf{Sample Output:}
\begin{verbatim}
Enter a positive integer: 5
Sum from 1 to 5 = 15
\end{verbatim}

%- - - - - - - - - - - - - - - - - - - - - - - - - - - - - - - - - - - - - -
\subsubsection{Trace Table for \texttt{sum(5)}}
%- - - - - - - - - - - - - - - - - - - - - - - - - - - - - - - - - - - - - -

\noindent\textbf{Phase 1 --- Recursive descent (pushing calls onto the stack):}

\begin{center}
\renewcommand{\arraystretch}{1.3}
\begin{tabular}{@{}clcl@{}}
\toprule
\textbf{Call} & \textbf{Function Call}
    & \textbf{Condition}
    & \textbf{Action} \\
\midrule
1 & \texttt{sum(5)} & $5 \neq 0$
    & returns $5 + \texttt{sum}(4)$ --- waits \\
2 & \texttt{sum(4)} & $4 \neq 0$
    & returns $4 + \texttt{sum}(3)$ --- waits \\
3 & \texttt{sum(3)} & $3 \neq 0$
    & returns $3 + \texttt{sum}(2)$ --- waits \\
4 & \texttt{sum(2)} & $2 \neq 0$
    & returns $2 + \texttt{sum}(1)$ --- waits \\
5 & \texttt{sum(1)} & $1 \neq 0$
    & returns $1 + \texttt{sum}(0)$ --- waits \\
6 & \texttt{sum(0)} & $0 = 0$
    & returns $0$ \quad (base case reached) \\
\bottomrule
\end{tabular}
\end{center}

\noindent\textbf{Phase 2 --- Stack unwinding (returning values):}

\begin{center}
\renewcommand{\arraystretch}{1.3}
\begin{tabular}{@{}clcr@{}}
\toprule
\textbf{Return} & \textbf{Returning From}
    & \textbf{Computation}
    & \textbf{Value} \\
\midrule
6 & \texttt{sum(0)} & base case
    & $0$ \\
5 & \texttt{sum(1)} & $1 + 0$
    & $1$ \\
4 & \texttt{sum(2)} & $2 + 1$
    & $3$ \\
3 & \texttt{sum(3)} & $3 + 3$
    & $6$ \\
2 & \texttt{sum(4)} & $4 + 6$
    & $10$ \\
1 & \texttt{sum(5)} & $5 + 10$
    & $\mathbf{15}$ \\
\bottomrule
\end{tabular}
\end{center}

\noindent\textbf{Verification using the closed-form formula:}
\[
    \texttt{sum}(5) = \frac{5 \times 6}{2} = \boxed{15}
\]

%-----------------------------------------------------------------------------
\subsection{Example: Recursive Factorial}
\label{subsec:recursive-factorial}
%-----------------------------------------------------------------------------

The factorial of a non-negative integer $n$ is defined as:
\[
    n! =
    \begin{cases}
        1 & \text{if } n = 0 \text{ or } n = 1
            \quad \text{(base case)} \\
        n \times (n-1)! & \text{if } n > 1
            \quad \text{(recursive case)}
    \end{cases}
\]

\begin{cppcode}
#include <iostream>
using namespace std;

long long factorial(int n)
{
    if (n <= 1)
        return 1;           // Base case: 0!=1!=1
    else
        return n * factorial(n - 1); // Recursive
}

int main()
{
    int n;
    cout << "Enter a non-negative integer: ";
    cin >> n;

    cout << n << "! = " << factorial(n) << endl;

    return 0;
}
\end{cppcode}

\noindent\textbf{Sample Output:}
\begin{verbatim}
Enter a non-negative integer: 6
6! = 720
\end{verbatim}

%- - - - - - - - - - - - - - - - - - - - - - - - - - - - - - - - - - - - - -
\subsubsection{Trace Table for \texttt{factorial(6)}}
%- - - - - - - - - - - - - - - - - - - - - - - - - - - - - - - - - - - - - -

\noindent\textbf{Phase 1 --- Recursive descent:}

\begin{center}
\renewcommand{\arraystretch}{1.3}
\begin{tabular}{@{}clcl@{}}
\toprule
\textbf{Call} & \textbf{Function Call}
    & \textbf{Condition}
    & \textbf{Action} \\
\midrule
1 & \texttt{factorial(6)} & $6 > 1$
    & returns $6 \times \texttt{factorial}(5)$ \\
2 & \texttt{factorial(5)} & $5 > 1$
    & returns $5 \times \texttt{factorial}(4)$ \\
3 & \texttt{factorial(4)} & $4 > 1$
    & returns $4 \times \texttt{factorial}(3)$ \\
4 & \texttt{factorial(3)} & $3 > 1$
    & returns $3 \times \texttt{factorial}(2)$ \\
5 & \texttt{factorial(2)} & $2 > 1$
    & returns $2 \times \texttt{factorial}(1)$ \\
6 & \texttt{factorial(1)} & $1 \leq 1$
    & returns $1$ \quad (base case) \\
\bottomrule
\end{tabular}
\end{center}

\noindent\textbf{Phase 2 --- Stack unwinding:}

\begin{center}
\renewcommand{\arraystretch}{1.3}
\begin{tabular}{@{}clcr@{}}
\toprule
\textbf{Return} & \textbf{Returning From}
    & \textbf{Computation}
    & \textbf{Value} \\
\midrule
6 & \texttt{factorial(1)} & base case
    & $1$ \\
5 & \texttt{factorial(2)} & $2 \times 1$
    & $2$ \\
4 & \texttt{factorial(3)} & $3 \times 2$
    & $6$ \\
3 & \texttt{factorial(4)} & $4 \times 6$
    & $24$ \\
2 & \texttt{factorial(5)} & $5 \times 24$
    & $120$ \\
1 & \texttt{factorial(6)} & $6 \times 120$
    & $\mathbf{720}$ \\
\bottomrule
\end{tabular}
\end{center}

\noindent\textbf{Verification:}
\[
    6! = 6 \times 5 \times 4 \times 3 \times 2 \times 1
       = \boxed{720}
\]

\begin{warningbox}{Stack Overflow Risk}
Each recursive call adds a new frame to the program's \textbf{call stack}.
For very large inputs (e.g., \texttt{factorial(100000)}), the stack may run
out of memory, causing a \emph{stack overflow} crash. Always ensure that:
\begin{itemize}
    \item The base case is reachable for all valid inputs.
    \item Each recursive call moves closer to the base case.
    \item The expected recursion depth is within practical limits.
\end{itemize}
\end{warningbox}

\begin{conceptbox}{When to Use Recursion}
Recursion is most natural for problems with a \emph{recursive structure}:
\begin{itemize}
    \item \textbf{Mathematical definitions:} Factorial, Fibonacci, power
          functions.
    \item \textbf{Data structures:} Traversing linked lists, binary trees,
          and graphs.
    \item \textbf{Divide-and-conquer algorithms:} Merge sort, quick sort,
          binary search.
\end{itemize}
For simple sequential tasks, an iterative solution (\texttt{for}/\texttt{while}
loop) is usually more efficient and easier to understand.
\end{conceptbox}

%-----------------------------------------------------------------------------
\section{Series Computation Using Functions}
\label{sec:series-computation}
%-----------------------------------------------------------------------------

A common programming exercise is to compute the value of a mathematical
series term by term. The following example computes the series:
\[
    S = x - \frac{x^{3}}{3!} + \frac{x^{5}}{5!}
      - \frac{x^{7}}{7!} + \cdots
      \pm \frac{x^{n}}{n!}
\]
where $n$ is an odd positive integer. This is related to the Taylor series
expansion of $\sin(x)$.

\begin{cppcode}
#include <iostream>
#include <cmath>
using namespace std;

// Compute factorial of n
long long fact(int n)
{
    long long result = 1;
    for (int i = 2; i <= n; i++)
        result *= i;
    return result;
}

// Compute x raised to the power p
double power(double x, int p)
{
    double result = 1.0;
    for (int i = 0; i < p; i++)
        result *= x;
    return result;
}

int main()
{
    double x;
    int n;

    cout << "Enter value of x (radians): ";
    cin >> x;
    cout << "Enter number of terms n: ";
    cin >> n;

    double sum = 0.0;
    int sign = 1;

    for (int i = 0; i < n; i++)
    {
        int exp = 2 * i + 1;  // 1, 3, 5, 7, ...
        sum += sign * power(x, exp) / fact(exp);
        sign *= -1;            // Alternate sign
    }

    cout << "Series sum  = " << sum << endl;
    cout << "sin(" << x << ") = " << sin(x) << endl;

    return 0;
}
\end{cppcode}

\noindent\textbf{Sample Output:}
\begin{verbatim}
Enter value of x (radians): 1.0
Enter number of terms n: 5
Series sum  = 0.841471
sin(1.0) = 0.841471
\end{verbatim}

\begin{infobox}{Understanding the Series Terms}
For $x = 1.0$ and 5 terms, the computation proceeds as follows:

\begin{center}
\renewcommand{\arraystretch}{1.3}
\begin{tabular}{@{}ccccr@{}}
\toprule
\textbf{Term} & \textbf{Exponent}
    & \textbf{Sign} & \textbf{Expression}
    & \textbf{Value} \\
\midrule
1 & 1 & $+$ & $\dfrac{1^{1}}{1!} = \dfrac{1}{1}$
    & $1.000000$ \\[6pt]
2 & 3 & $-$ & $\dfrac{1^{3}}{3!} = \dfrac{1}{6}$
    & $-0.166667$ \\[6pt]
3 & 5 & $+$ & $\dfrac{1^{5}}{5!} = \dfrac{1}{120}$
    & $0.008333$ \\[6pt]
4 & 7 & $-$ & $\dfrac{1^{7}}{7!} = \dfrac{1}{5040}$
    & $-0.000198$ \\[6pt]
5 & 9 & $+$ & $\dfrac{1^{9}}{9!} = \dfrac{1}{362880}$
    & $0.000003$ \\
\midrule
  &   &     & \textbf{Sum}
    & $\mathbf{0.841471}$ \\
\bottomrule
\end{tabular}
\end{center}
\end{infobox}

%-----------------------------------------------------------------------------
\section{Chapter Summary}
\label{sec:ch2-summary}
%-----------------------------------------------------------------------------

\begin{importantbox}{Key Takeaways -- Advanced Function Concepts}
\begin{enumerate}
    \item User-defined functions fall into \textbf{three categories}:
          no args / no return, args / no return, and args / with return.
    \item \textbf{Actual arguments} are the values passed in a function call;
          \textbf{formal arguments} are the parameters in the function
          definition that receive those values.
    \item \textbf{Local variables} exist only within their enclosing block;
          \textbf{global variables} are accessible throughout the program.
    \item \textbf{Recursive functions} call themselves and must have both
          a \emph{base case} and a \emph{recursive case}.
    \item Always trace recursive calls using a \textbf{descent / unwind table}
          to verify correctness and understand the call stack.
    \item Functions enable \textbf{series computation} by encapsulating
          repetitive mathematical operations (factorial, power) into
          reusable units.
\end{enumerate}
\end{importantbox}
